\newpage
\section{问题分析}

\subsection{问题一：交通图构建分析}
本题的核心是将城市交通网络抽象为图结构G=(V,E)。
交叉口作为节点v_i，路段作为边e_ij。
节点特征包括历史流量、时间信息、天气等；
边特征包括路段长度、限速、实时流量等。
关键挑战在于如何处理大规模图的表示学习。

\subsection{问题二：短时预测建模}
短时交通流预测需要捕捉时空相关性。
GraphSAGE通过聚合邻居信息学习节点嵌入：
h_v^{(l+1)} = sigma(W^{(l)} · CONCAT(h_v^{(l)}, AVG({h_u^{(l)}}_{u in N(v)})))
该聚合方式支持归纳式学习，可处理新节点。

\subsection{问题三：模型验证策略}
采用时间序列交叉验证，将数据按时间划分训练集/验证集/测试集。
评估指标选用MAE、RMSE、MAPE。
与ARIMA、LSTM、GCN等基线模型对比。
